\scanpage{208}
\uline{Critère de type fini}
\margin{$\deleted{\unsure{\beta}}\unsure{\mu|\beta}(s)+2\alpha(s)$ \unsure{$1/2$} continue}

\uuline{Proposition} \quad Soit $f:X\longrightarrow Y$ un morphisme loc. de
type fini de préschémas loc noethériens. On
suppose (i) $f$ séparé (ii) $f$ \deleted{\unsure{univers.}} ouvert (iii) les
fibres de $f$ sont \deleted{\unsure{de type fini}}\inserted{propres sur les $k(y)$} Alors $f$
est de type fini.
% slanted
\margin{prouver que $\underline{\mathrm{Pic}}^{\unsure{\tau}}_{X/S}$ est quasi-compact si $f:X\unsure{-|\to}S$ est \unsure{comment|comme ça} ?}

\inserted{On peut supposer $Y$ affine, donc \unsure{noeth.}} Il suffit de prouver $f$ qu. cpct, i.e $X$ quasi cpct.
Une \unsure{récurrence|récurrance} noeth. \deleted{\unsure{nous}} \unsure{$1/$} nous
ramène à prouver ceci : \deleted{il existe si $Y$ \unsure{en}}\inserted{Il existe}
un ouvert \unsure{non} vide de $Y$ tel que $X|U$ soit
qu. cpct. \inserted{On peut supposer $Y$ intègre, pt gén. $\eta$} Il existe un ouvert qu. cpct $V$ de $X$ \deleted{de
type fini} contenant $X_\eta$. Il existe un
$X'$ \deleted{de type fini}\inserted{\uline{propre}} sur $Y$, et un \unsure{$k(\eta)$-}isom
$X'_\eta\simeq X_\eta=V_\eta$. Si on restreint $Y$, on peut
supposer que le dernier isom provient
d'un isomorphis $X'\simeq V$. Donc \deleted{$V$ \unsure{soit} propre}
$V$ est propre sur $Y$, donc fermé dans $X$ [\unsure{car|où} $X$ est
séparé sur $\deleted{\unsure{X}}Y$]. \deleted{\unsure{On a un pt}} \deleted{Alors $X$\unsure{$-V$|$\,{}_\eta$} est une partie}
\deleted{\unsure{f} donc} $X-V$ est \uline{ouvert}. \deleted{Comme $X$}\inserted{Comme $X\unsure{-|\to}Y$}
\deleted{\unsure{la} \unsure{lettre|loi} dernière \unsure{donne}} est ouvert, \deleted{$X-V$ donc}
si on avait $X-V\ne\varnothing$, $X-V$ dominerait $Y$,
\deleted{donc} ce qui est incompatible avec $X_\eta=V_\eta$
i.e. $(X-V)_\eta=\varnothing$. Donc $X=V$, d'où la conclusion
% archive sheet number: 207
